\SourceLeaf{462}nicht ohne Weiteres dazu dienen, uns zu belehren, ob und wie viel die Läger und Pfannen noch verrückt werden müssen, um den Bedingungen 3 und 6 im 3. Art. Genüge zu leisten, indem diese sich auf den Mittelpunkt des innern Kreises, jene aber auf den Mittelpunkt der Eintheilung beziehen, welche beide etwas verschieden sein können, und an dem in Rede stehenden Instrumente auch wirklich verschieden sind. In der That waren vor den hier angeführten Beobachtungen die betreffenden Berichtigungen mit aller möglichen Sorgfalt ausgeführt.

Die mit $z$ bezeichnete Grösse ist offenbar für jede Nadel unveränderlich, und eine ähnliche Behandlung der Beobachtungen von andern Tagen hat nahe denselben Werth ergeben Für die drei andern Nadeln habe ich gefunden
\begin{center}
\begin{tabular}{crll}
für Nadel&$2$&$.\quad.\quad.\quad.\quad.\quad.\quad.\quad.$&$+3'\ 18''$\\
&$3$&$.\quad.\quad.\quad.\quad.\quad.\quad.\quad.$&$-1\ 4$\\
&$4$&$.\quad.\quad.\quad.\quad.\quad.\quad.\quad.$&$+1\ 2$
\end{tabular}
\end{center}

Obwohl die Kenntniss dieser Werthe kein besonderes praktisches Interesse hat, so gibt doch ihre Kleinheit ein rühmliches Zeugniss für die von dem ausgezeichneten Künstler auf die Bearbeitung der Nadeln verwandte Sorgfalt.

\Article{11}
Das Mittel der Ablesungen der beiden Spitzen gibt uns die Neigung der diese Spitzen verbindenden geraden Linie oder einer Parallele mit derselben gegen den mit $0$ bezeichneten Diameter des Verticalkreises. Ich stelle diese $16$ Mittel hier paarweise zusammen.
\begin{center}
\textit{Nadelende $B$ Nordpol}\par\smallskip
{\fontsize{9}{11}\selectfont
\begin{tabular}{c|c||c|c}
Azim.&Bez. Nadelfl.&Azim.&Bez. Nadelfl.\\
Kr.&vorne&Kr.&hinten\\\hline
$90^0\ 5'$&$67^0\ 28'\ 45''$&$270^0\ 5'$&$112^0\ 36'\ 23''$\\
$180\ 5$&$89\ 52\ 45$&$0\ 5$&$90\ 12\ 11$\\
$270\ 5$&$112\ 17\ 43$&$90\ 5$&$67\ 45\ 15$\\
$0\ 5$&$89\ 58\ 10$&$180\ 5$&$90\ 11\ 30$
\end{tabular}}
\par\medskip
\textit{Nadelende $A$ Nordpol}\par\smallskip
{\fontsize{9}{11}\selectfont
\begin{tabular}{c|c||c|c}
$90^0\ 5'$&$68^0\ 2'\ 42''$&$270^0\ 5'$&$112^0\ 6'\ 20''$\\
$180\ 5$&$90\ 13\ 34$&$0\ 5$&$89\ 54\ 6$\\
$270\ 5$&$112\ 24\ 57$&$90\ 5$&$67\ 36\ 7$\\
$0\ 5$&$90\ 15\ 8$&$180\ 5$&$89\ 51\ 24$
\end{tabular}}
\end{center}

\SourceLeaf{463}Nebeneinander stehen hier diejenigen Einstellungen, bei welchen in entgegengesetzter Lage des Verticalkreises die Zapfenachse gleiche Lage (gegen die Weltgegenden) hatte. Der Zusammenhang zweier solcher Zahlen $l$ und $l'$ ist ein sehr einfacher, wenn die Läger so berichtigt sind, dass eine gegen die verticale Drehungsachse normale Ebene sie berührt. In dieser Voraussetzung liegt in beiden Einstellungen die Zapfenachse auf einer horizontalen Ebene und der Ruhestand der Nadel ist daher offenbar derselbe, d. i. wenn wir unter $L$ die Neigung der von der obern zur untern Spitze gezogenen geraden Linie gegen denjenigen horizontalen Radius des Kreises verstehen, der jedesmal auf der rechten Seite der gezeichneten Nadelfläche liegt, so wird $L$ in beiden Einstellungen gleiche Werthe haben. Diese Neigung ergibt sich aber
\[
\begin{aligned}
\text{in der ersten Einstellung }&=l-\alpha\\
\text{und in der zweiten }&=180^0-(l'-\alpha)
\end{aligned}
\]
wenn $\alpha$ den Fehler des Nullpunkts (d. i. die Ablesung an demjenigen Kreisradius, der mit der Verticalachse einen rechten Winkel macht) bedeutet. Wir haben also unter obiger Voraussetzung
\[
\alpha=\tfrac{1}{2}(l+l')-90^0,\qquad L=\tfrac{1}{2}(l+180^0-l')
\]
Aus den Beobachtungen vom 23. Sept., wo diese Berichtigung mit Hülfe der im 5. Art. erwähnten Vorrichtung auf das sorgfältigste ausgeführt war, erhalten wir also acht verschiedene Bestimmungen von $\alpha$, nemlich
\begin{center}
\begin{tabular}{r}
$+2'\ 34''$\\
$2\ 28$\\
$1\ 29$\\
$4\ 50$\\
$4\ 31$\\
$3\ 50$\\
$0\ 32$\\
$3\ 16$
\end{tabular}
\end{center}

Die Summe der Quadrate der in Secunden ausgedrückten Abweichungen von dem Mittelwerthe $2'56''$ findet sich $=57214$; wenn man also diese Abweichungen wie ganz zufällige betrachtet, so ergeben sie den mittlern Fehler des
\SourceLeaf{464}Resultats aus einem Paar coordinirten Einstellungen $=\BareRoot{\frac{57214}{7}}=90''4$. Man sieht, dass bei diesem Instrumente die Anomalien der Einstellung viel beträchtlicher sind, als die reinen Ablesungsfehler.

\Article{12}
Anders verhält es sich aber, wenn die vorausgesetzte genaue Berichtigung der Läger nicht Statt findet. Nehmen wir an, dass zwar die Ränder derselben in Einer Ebene liegen, aber nicht in einer gegen die Verticalachse normalen, so ist in den beiden Einstellungen diese Ebene auf entgegengesetzte Art gegen die Horizontalebene geneigt. Hier kommt indessen nur die Neigung in dem Sinn der Lagerränder oder parallel mit der Kreisebene in Betracht, indem eine kleine Neigung in der Querrichtung oder in dem Sinn der Nadelachse keinen merklichen Einfluss auf die Ruhestellung der Nadel hat. Es bezeichne nun $L$ diejenige Neigung der Nadel (eben so verstanden wie oben), welche bei dem Aufliegen auf einem vollkommen horizontalen Lager Statt finden würde; $\delta$ die entsprechende Richtungskraft, d. i. den Coëfficienten, in welchen der Sinus einer Ablenkung von der Ruhestellung multiplicirt werden muss, um das Drehungsmoment der die Nadel nach dieser Stellung zurücktreibenden Kraft auszudrücken; endlich sei $L+\beta$ die in der ersten Einstellung auf dem geneigten Lager wirklich Statt findende Neigung. Es lässt sich dann leicht zeigen, dass
\[
\delta\sin\beta=p\rho\sin\gamma
\]
wird, wo $p$ das Gewicht der Nadel, $\rho$ den Halbmesser der Zapfen und $\gamma$ die Neigung des Lagers gegen die Horizontallinie bedeuten, letztere Grösse positiv genommen, wenn das Lager auf der rechten Seite der gezeichneten Nadelfläche niedriger ist. Offenbar muss nun aber in der zweiten Einstellung $-\gamma$ anstatt $\gamma$ gesetzt werden, wodurch $\beta$ in $-\beta$ übergeht, daher in dieser zweiten Einstellung die Neigung der Nadel $L-\beta$ wird. Wir haben also
\[
l-\alpha=L+\beta,\qquad180^0-(l'-\alpha)=L-\beta
\]
und folglich, eben so wie im vorhergehenden Art.
\[
\tfrac{1}{2}(l+180^0-l')=L
\]
hingegen anstatt der andern dortigen Gleichung jetzt
\[
\tfrac{1}{2}(l+l')-90^0=\alpha+\beta
\]


\expandafter\def\csname HistoricalNoteText465\endcsname{Eigentlich ist es die Summe der Trägheitsmomente der Nadel und des Bügels; beide von einander zu scheiden ist theils unthunlich, theils überflüssig, da keine andere Schwingungen als horizontale mit diesem Bügel gebraucht werden.}
\SourceLeaf{465}Liegen aber die Ränder der Achatplatten gar nicht in Einer Ebene, so werden eben diese beiden Formeln auch noch hinlänglich genau gültig bleiben, wenn man nur für $\gamma$ das Mittel der Neigungen der beiden Kanten annimmt, vorausgesetzt, dass der Schwerpunkt der Nadel von den beiden aufliegenden Punkten der Zapfen nahe gleich weit absteht. Genau genommen entsteht zwar noch eine kleine Modification aus dem Umstande, dass dann die gerade Linie, welche die beiden Berührungspunkte der Zapfen und Läger verbindet, in den beiden Einstellungen nicht ganz gleiche Azimuthe hat; der Einfluss dieses Umstandes auf die Stellung der Nadel wird aber auch da, wo er am stärksten ist, nemlich bei Beobachtungen in der gegen den magnetischen Meridian normalen Ebene, wie ganz unmerklich betrachtet werden dürfen.

\Article{13}
Da es nicht uninteressant ist, übersehen zu können, in welchem Verhältnisse bei nicht berichtigtem Zustande der Läger die Neigung derselben auf die Einstellung der Nadel wirkt, so füge ich hier noch das dazu nöthige für die am 23. Sept. gebrauchte Nadel bei. Zu dem Zweck, ihr Trägheitsmoment zu bestimmen, hatte ich schon früher horizontale Schwingungen derselben beobachtet, theils ohne, theils mit Auflegung eines Ringes, dessen eignes Trägheitsmoment sich aus Gewicht und Dimensionen mit hinlänglicher Schärfe berechnen liess. Es war am 21. September
\begin{center}
{\fontsize{9}{11}\selectfont
\begin{tabular}{llrl}
Schwingungsdauer ohne Ring&$.\quad.\quad.\quad.$&$5''88431$&\\
\hspace{2em}---\quad mit Ring&$.\quad.\quad.\quad.$&$7.32835$&\\
Gewicht des Ringes&$.\quad.\quad.\quad.\quad.\quad.\quad.$&$19.2385$&Gramme\\
Innerer Durchmesser&$.\quad.\quad.\quad.\quad.\quad.$&$75.525$&Millimeter\\
Äusserer Durchmesser&$.\quad.\quad.\quad.\quad.\quad.$&$79.767$&---
\end{tabular}}
\end{center}

\Needspace{6\baselineskip}
Hieraus folgt, Gramm und Millimeter als Einheit angenommen,
\begin{center}
\begin{tabular}{llr}
Trägheitsmoment des Ringes&$.\quad.\quad.\quad.$&$29019$\\
\hspace{2em}---\quad der Nadel\footnotemark&$.\quad.\quad.\quad.$&$52662$
\end{tabular}
\end{center}
\footnotetext{\csname HistoricalNoteText465\endcsname}


\SourceLeaf{466}Hieraus verbunden mit den oben Art. 9 angegebenen Schwingungszeiten vom 23. September, und die Länge des einfachen Secundenpendels in Göttingen zu $994.126$ Millimeter angenommen, ergibt sich auf bekannte Weise
\begin{center}
\begin{tabular}{llr}
\multicolumn{3}{l}{horizontale magnetische Richtungskraft}\\
\quad vor dem Umstreichen&$.\quad.\quad.\quad.$&$1.5556$\\
\quad nach dem Umstreichen&$.\quad.\quad.\quad.$&$1.5352$
\end{tabular}
\end{center}

Diese Zahlen gelten, genau genommen. zunächst nur für den Platz, wo die Schwingungen beobachtet sind, und schliessen also die daselbst etwa statt findenden localen Einflüsse ein: für den gegenwärtigen Zweck kommt dieser jedenfalls nur geringe Einfluss nicht in Betracht.

Mit Neigung $67^0\ 40'\ 54''$ folgt hieraus ferner
\begin{center}
\begin{tabular}{llr}
\multicolumn{3}{l}{ganze magnetische Richtungskraft,}\\
\quad vor dem Umstreichen&$.\quad.\quad.\quad.$&$4.0965$\\
\quad nach dem Umstreichen&$.\quad.\quad.\quad.$&$4.0429$\\
\multicolumn{3}{l}{verticale magnetische Richtungskraft}\\
\quad vor dem Umstreichen&$.\quad.\quad.\quad.$&$3.7897$\\
\quad nach dem Umstreichen&$.\quad.\quad.\quad.$&$3.7401$
\end{tabular}
\end{center}

Diese vier Zahlen können, wenn man die kleine Modification, welche die magnetische Richtungskraft der Nadel durch die Excentricität des Schwerpunkts erhält, nicht berücksichtigt, als die Werthe von $\delta$ betrachtet werden, je nachdem die Beobachtung im magnetischen Meridian oder in der dagegen normalen Ebene gemacht ist. Da $\beta$ und $\gamma$ immer klein genug sind, um diese Grössen selbst an die Stelle ihrer Sinus setzen zu können, also
\[
\beta=\frac{p\rho}{\delta}.\gamma
\]
so ergibt sich hieraus, je nachdem die Stärke der Magnetisirung, wie sie vor oder wie sie nach dem Umstreichen war, zum Grunde gelegt wird

für Beobachtungen im magnetischen Meridian
\[
\beta=1{,}1882\gamma\quad\text{oder}\quad\beta=1{,}2039\gamma
\]
für Beobachtungen in der gegen den magnetischen Meridian normalen Ebene
\[
\beta=1{,}2844\gamma\quad\text{oder}\quad\beta=1{,}3014\gamma
\]


\SourceLeaf{467}Übrigens sind zwar die bisher betrachteten Relationen zwischen den einzelnen Beobachtungsstücken nicht wesentlich, insofern es nur gilt, aus allen die magnetische Inclination abzuleiten: allein sie sind nicht unwichtig für die Prüfung und Befestigung des Resultats, indem das rechte Vertrauen in das Ganze erst aus der klaren Einsicht in die befriedigende Übereinstimmung der Theile erwachsen kann.

\Article{14}
Die Ausbeute der Beobachtungen ist nunmehr auf die acht Werthe von $L$ zurückgeführt, welche erklärt werden können als die Neigungen der von der Südpolspitze der Nadel nach der Nordpolspitze gezogenen geraden Linie gegen den auf der rechten Seite der gezeichneten Nadelfläche liegenden horizontalen Kreisradius im Zustande des Gleichgewichts, insofern die Nadelzapfen auf einer horizontalen Fläche aufliegend gedacht werden, oder, was in statischer Rücksicht offenbar ganz dasselbe ist, insofern die Nadel als nur um die Achsenlinie der Zapfen drehbar angenommen wird. Mit andern Worten, die Werthe von $L$ sind die verbesserten d. i. vom Einfluss des Fehlers des Nullpunkts und der Nichthorizontalität der Läger befreieten Werthe der im 11. Art. unter der Überschrift \textit{Bezeichnete Nadelfläche vorne} aufgeführten Zahlen
\begin{center}
\textit{Werthe von $L$.}\par\smallskip
\begin{tabular}{c|c|c}
Az. Kr.&$B$ Nordpol&$A$ Nordpol\\\hline
$90^0\ 5'$&$67^0\ 26'\ 11''$&$67^0\ 58'\ 11''$\\
$180\ 5$&$89\ 50\ 17$&$90\ 9\ 44$\\
$270\ 5$&$112\ 16\ 14$&$112\ 24\ 25$\\
$0\ 5$&$89\ 53\ 20$&$90\ 11\ 52$
\end{tabular}
\end{center}

Um nun den Zusammenhang der Werthe von $L$ mit den Elementen, von welchen er abhängt, in einer Gleichung auszudrücken, bediene ich mich folgender Bezeichnungen.

$V$\quad Stellung des Azimuthalkreises für die Beobachtung.

$V^0$\quad Stellung des Azimuthalkreises, bei welcher der Verticalkreis im magnetischen Meridian, und die getheilte Seite nach Osten gerichtet ist.

$i$\quad magnetische Inclination.


\SourceLeaf{468}$m$\quad das Product des magnetischen Moments der Nadel in die ganze Intensität der erdmagnetischen Kraft, wobei die Schwere als Einheit der beschleunigenden Kräfte angenommen wird.

$q$\quad das Gewicht der Nadel multiplicirt in die Entfernung des Schwerpunkts von der Zapfenachse.

$c$\quad der spitze Winkel zwischen der die Spitzen der Nadel verbindenden geraden Linie und der magnetischen Achse derselben, positiv, wenn letztere rechts liegt, indem die Nadel mit der gezeichneten Seite nach oben horizontal liegend gedacht wird.

$Q$\quad der Winkel zwischen der geraden Linie von der Südpolspitze der Nadel nach der Nordpolspitze einerseits und der geraden Linie von der Zapfenachse nach dem Schwerpunkt andererseits, so verstanden, dass man von der ersten anfangend bei derselben Lage der Nadel wie für $c$ von der Linken nach der Rechten zählt.

$\delta$\quad die Richtungskraft.

Zerlegt man die erdmagnetische Kraft in einen verticalen und einen horizontalen Theil, so entsteht aus dem erstern das Drehungsmoment, positiv genommen in dem Sinn wachsender $L$,
\[
m\sin i\cos(L+c)
\]
aus dem andern
\[
-m\cos i\cos(V-V^0)\sin(L+c)
\]
Die Schwere hingegen bewirkt das Drehungsmoment
\[
q\cos(L+Q)
\]
Da $L$ die Gleichgewichtstellung ausdrückt, so wird die Summe dieser drei Momente $=0$; woraus wir die Hauptgleichung erhalten
\[
-\sin i\cos(L+c)+\cos i\cos(V-V^0)\sin(L+c)=\frac{q}{m}.\cos(L+Q)
\]
Schreiben wir in der Summe der drei Momente $L+z$ anstatt $L$, so erhalten wir das Drehungsmoment, welches bei einer Ablenkung $z$ von der Gleichgewichtstellung Statt findet; entwickelt man diesen Ausdruck in zwei Theile mit den Factoren $\cos z$ und $\sin z$, so verschwindet der erste vermöge der Hauptgleichung, und der zweite wird in Folge des Begriffs der Richtungskraft $=-\delta\sin z$.
\SourceLeaf{469}Wir haben also für $\delta$ die allgemeine Formel
\[
\delta=m\sin i\sin(L+c)+m\cos i\cos(V-V^0)\cos(L+c)+q\sin(L+Q)
\]
Für die drei speciellen Hauptfälle finden wir hieraus:

I.\quad Für $V=V^0$
\[
\begin{aligned}
\sin(L+c-i)&=\frac{q}{m}\cos(L+Q)\\
\delta&=m\cos(L+c-i)+q\sin(L+Q)\\
&=\frac{m\cos(Q+i-c)}{\cos(L+Q)}=\frac{q\cos(Q+i-c)}{\sin(L+c-i)}
\end{aligned}
\]

II.\quad Für $V=V^0+180^0$
\[
\begin{aligned}
\sin(L+c+i)&=-\frac{q}{m}\cos(L+Q)\\
\delta&=-m\cos(L+c+i)+q\sin(L+Q)\\
&=-\frac{m\cos(Q-c-i)}{\cos(L+Q)}=\frac{q\cos(Q-c-i)}{\sin(L+c+i)}
\end{aligned}
\]

III.\quad Übereinstimmend für $V=V^0+90^0$ und $V=V^0+270^0$
\[
\begin{aligned}
\sin i\cos(L+c)&=-\frac{q}{m}\cos(L+Q)\\
\delta&=m\sin i\sin(L+c)+q\sin(L+Q)\\
&=-\frac{m\sin i\sin(Q-c)}{\cos(L+Q)}=\frac{q\sin(Q-c)}{\cos(L+c)}
\end{aligned}
\]

Unser Beispiel gibt für die beiden letzten Fälle anstatt gleicher Werthe von $L$ Ungleichheiten von resp. $3'\ 3''$ und $2'\ 8''$, welche theils in den zufälligen Beobachtungsfehlern, theils in der Conspiration mehrerer Umstände ihren Grund haben: in einer kleinen Unsicherheit der anfänglichen magnetischen Orientirung; in der Veränderlichkeit der magnetischen Declination und also des Werthes von $V^0$ im Laufe der Beobachtungen; in einer kleinen Excentricität des Horizontalkreises, welche in Ermangelung einer doppelten Ablesung nicht controllirt werden kann; endlich darin, dass die Rechtwinkligkeit der Zapfenachse gegen die Kreisebene durch die Auflegung vermittelst der Pfannen nur auf eine unvollkommene Art erhalten werden kann. Alle diese Umstände werden, so viel thunlich, unschädlich gemacht, indem man aus beiden Einstellungen die Mittel nimmt, also

\SourceLeaf{470}\begin{center}
\begin{tabular}{llr}
für $B$ Nordpol&$.\quad.\quad.\quad.\quad.$&$L=89^0\ 51'\ 49''$\\
für $A$ Nordpol&$.\quad.\quad.\quad.\quad.$&$L=90\ 10\ 48$
\end{tabular}
\end{center}
setzt. Indessen wird man dieser Umstände wegen immer dem Resultate für die Einstellung bei einer gegen den magnetischen Meridian normalen Lage eine etwas geringere Zuverlässigkeit beilegen müssen, als bei den Lagen im Meridian selbst, wo der Einfluss jener Ursachen als unmerklich betrachtet werden kann.

\Article{15}
Die aus den $32$ ursprünglichen Zahlen uns übrig gebliebenen sechs mögen fortan auf folgende Art bezeichnet werden:
\begin{center}
\begin{tabular}{l|l}
Werthe von $L$&für $V-V^0=$\\\hline
$f,\ f'$&$0$\\
$180^0-g,\ 180^0-g'$&$180^0$\\
$h,\ h'$&$90^0$ und $270^0$
\end{tabular}
\end{center}
wo die nicht accentuirten Zeichen sich auf $B$ Nordpol, die accentuirten auf $A$ Nordpol beziehen sollen. Offenbar sind so $f,f',g,g'$ für die Stellungen im magnetischen Meridian die Neigungen der von der Südpolspitze der Nadel nach der Nordpolspitze gezogenen geraden Linie sämmtlich unter der nordlichen Horizontallinie, und zwar die beiden ersten für die Stellung, wo die gezeichnete Nadelfläche nach Osten gekehrt ist, die beiden andern für die entgegengesetzte; $h,h'$ hingegen sind, für die Stellungen in der gegen den magnetischen Meridian normalen Ebene, die Neigungen derselben geraden Linie gegen die östliche oder westliche Horizontallinie, je nachdem die gezeichnete Nadelfläche nach Süden oder nach Norden gekehrt ist.

Was die Elemente betrifft, von welchen diese sechs Grössen abhängen, so ist $q$ ganz constant, und $i$ muss für alle als gleich angenommen werden, insofern wir die im Laufe der Beobachtungen etwa Statt habenden kleinen Schwankungen doch nicht berücksichtigen können; $Q,m,c$ hingegen ändern nach dem Umstreichen ihre Werthe, und zwar $Q$ genau um $180^0$, $m$ und $c$ aber so, dass weiter kein bestimmter Zusammenhang mit den frühern Statt findet, als dass wir wenn zum Umkehren der Pole eine gleichförmige Streichmanipulation und kräf\SourceLeaf{471}tige Streichstäbe angewandt werden, versichert sein dürfen, dass der Unterschied und für $c$ auch die absoluten Werthe nicht sehr beträchtlich sein können. Indem ich nun fortan die nicht accentuirten Zeichen $Q,m,c$ die bestimmten für die Beobachtungen mit $B$ Nordpol geltenden Werthe bedeuten, und für die Beobachtung mit $A$ Nordpol, $Q+180^0,m',c'$ an ihre Stelle treten lasse, verwandeln sich die allgemeinen Gleichungen des vorhergehenden Art. in folgende sechs:
\[
\begin{array}{lll}
\sin(f+c-i)=\frac{q}{m}\cos(f+Q)&.\ .\ .\ .\ .\ .\ .&(1)\\[4pt]
\sin(g-c-i)=\frac{q}{m}\cos(g-Q)&.\ .\ .\ .\ .\ .\ .&(2)\\[4pt]
\sin i\cos(h+c)=-\frac{q}{m}\cos(h+Q)&\phantom{.}\ .\ .\ .\ .\ .\ .&(3)\\[4pt]
\sin(f'+c'-i)=-\frac{q}{m'}\cos(f'+Q)&\phantom{.}\ .\ .\ .\ .\ .\ .&(4)\\[4pt]
\sin(g'-c'-i)=-\frac{q}{m'}\cos(g'-Q)&\phantom{.}\ .\ .\ .\ .\ .\ .&(5)\\[4pt]
\sin i\cos(h'+c')=\frac{q}{m'}\cos(h'+Q)&.\ .\ .\ .\ .\ .\ .\ .&(6)
\end{array}
\]

\Article{16}
Theoretisch betrachtet reichen diese sechs Gleichungen hin, um die sechs unbekannten Grössen $c,c',\frac{q}{m},\frac{q}{m'},Q,i$ zu bestimmen, und es mag der Auflösung dieser Aufgabe ein Platz hier vergönnt sein, obgleich sie gar keinen praktischen Werth hat, da der enorme Einfluss der unvermeidlichen Beobachtungsfehler auf die Endresultate dieses Verfahren ganz unbrauchbar macht.

Multiplicirt man die Gleichungen $1,2,3$ resp. mit
\[
\sin(g+h),\quad\sin(f-h),\quad\sin(f+g)
\]
und addirt, so erhält man nach einigen leichten Reductionen
\[
\sin(f+c).\sin(g+h)=\sin(g-c).\sin(h-f)
\]
woraus sich $c$ leicht bestimmen lässt, am bequemsten vermittelst der Formel
\[
\operatorname{tang}(c+\tfrac{1}{2}f-\tfrac{1}{2}g)=-\operatorname{tang}\tfrac{1}{2}(f+g)^2.\operatorname{cotang}(h-\tfrac{1}{2}f+\tfrac{1}{2}g)
\]
Auf ähnliche Art erhält man aus den Gleichungen $4,5,6$
\[
\operatorname{tang}(c'+\tfrac{1}{2}f'-\tfrac{1}{2}g')=-\operatorname{tang}\tfrac{1}{2}(f'+g')^2.\operatorname{cotang}(h'-\tfrac{1}{2}f'+\tfrac{1}{2}g')
\]
